\chapter{Internals}

This chapter opens the implementation. You do not need any of it to write wick
games, but a language's design is best understood through its machinery, and
wick's machinery is small enough to explain honestly in a chapter. The whole
language --- lexer, compiler, virtual machine, garbage collector, built-ins ---
is about two thousand lines of dependency-free \mbox{C\texttt{++}} in three
files:

\begin{center}
\begin{tabular}{@{}ll@{}}
\toprule
\texttt{wick/wick.hpp} & the embed API (about 120 lines) \\
\texttt{wick/wick\_front.cpp} & lexer and one-pass typed compiler $\rightarrow$ bytecode \\
\texttt{wick/wick\_vm.cpp} & stack VM, mark/sweep GC, built-ins, signature parser \\
\bottomrule
\end{tabular}
\end{center}

\section{One-pass typed compilation}

The defining choice is that \textbf{there is no abstract syntax tree}. The
compiler is single-pass, in the style of Crafting Interpreters' \emph{clox}: a
Pratt expression parser emits bytecode as it goes, and it carries a static
`Type` for every expression as it emits. The type checker and the code
generator are not two passes over a tree; they are the same walk. When the
parser finishes an expression, it has both produced the bytecode that computes
the value and determined the value's type, together, in one sweep.

Two visible language rules fall directly out of this design.

\paragraph{Declare before use.} Because the compiler resolves each name and
each call at the moment it reaches it --- there is no later pass to fix up
forward references --- a function must be defined textually before it is
called, and a variable must be declared before it is used. This is not a
limitation the designers chose for its own sake; it is what single-pass
compilation costs, and the price is small.

\paragraph{Types are known at emit time.} Since every expression's type is in
hand as its code is generated, the compiler can \emph{specialise} the
bytecode. This is the performance story, and it is worth being concrete about.

\section{Specialised opcodes}

In a dynamically typed VM, `a + b` compiles to a single generic ``add''
opcode that, at run time, inspects the tags of its operands and decides
whether to add numbers, concatenate strings, or raise an error. That tag check
happens on every execution of the instruction. wick pays it \emph{once}, at
compile time, because it already knows the types:

\begin{itemize}
\item `+` on two `num` emits `OP_ADD`; `+` on two `str` emits `OP_CONCAT`.
There is never a runtime tag check to choose between them --- the choice was
made when the code was generated.
\item The ordered comparisons compile to numeric compare opcodes
(`OP_LT`, `OP_LE`, `OP_GT`, `OP_GE`), because their operands are known to be
`num`.
\item Local variables resolve to \emph{frame slots} --- small integer offsets
into the current call's stack frame --- so `OP_LOADL` and `OP_STOREL` take a
slot number, and there is no name lookup at run time. Globals resolve to slot
indices too, via `OP_LOADG` and `OP_STOREG`.
\item Engine natives resolve to integer ids at compile time. The VM's native
call, `OP_NCALL`, is an array index into the natives table, not a hash-table
walk by name.
\end{itemize}

\noindent
The bytecode has roughly forty opcodes. Alongside the arithmetic and
comparison ops there are the control-flow jumps (`OP_JMP`, `OP_JF` which pops a
`bool`, `OP_JB` for backward jumps in loops), two jumps that implement
optionals (`OP_JNN`, ``jump if not nil'', and `OP_ISNIL`), the call ops
(`OP_CALL` for wick functions, `OP_NCALL` for natives), the container ops
(`OP_LIST`, `OP_MAP`, and the index get/set pairs `OP_IGET`/`OP_ISET` and
`OP_MGET`/`OP_MSET`), and the specialised built-ins (`OP_LEN`, `OP_TOSTR`,
`OP_TONUM`, `OP_LPUSH`, `OP_LPOP`). Each corresponds to something the
type checker proved about the code around it.

\section{The stack VM}

The runtime is a classic stack machine. There is a constants pool, a stack of
values, and a stack of call frames; each frame records its base --- the point
in the value stack where its slots begin --- and its instruction pointer.
Arguments to a call are already on the stack when the frame is pushed, and the
callee's locals live in the slots just above them. Returning resizes the stack
back to the frame's base, so the stack is empty between top-level statements
and, crucially, empty between frames.

Runtime errors carry a bytecode-offset-to-line table, built during
compilation, so that when the VM raises an error --- an out-of-range index, a
failed `check` --- it can report the source line, and the engine's error screen
can point at it. This is why a wick runtime error reads `file:line: message`
and not a raw bytecode offset.

There is \textbf{no JIT}, and this is deliberate, not unfinished. The
interpreter is fast enough for game logic at 400 by 240 --- the work per frame
is small, and the specialised opcodes remove the overhead a naive interpreter
would carry. More importantly, the engine targets platforms, iOS chief among
them, where generating and executing code at run time is forbidden. A JIT
would make wick illegal on those platforms; the specialised-bytecode interpreter
is legal everywhere by construction, and its performance is predictable, with
no warm-up and no deoptimisation cliffs.

\section{Frame-boundary garbage collection}

wick manages three kinds of heap object --- strings, lists, and maps --- with a
mark-and-sweep collector. What makes it a game collector rather than a general
one is \emph{when} it runs: never on its own initiative, only when the host
calls `wick::collect()`. lantern calls that function exactly once per frame,
after presenting the image.

The consequence is that a collection can never land in the middle of a draw.
The one moment a pause is invisible --- after the frame is on screen, before the
next one begins --- is the only moment collection happens. The worst-case cost
of a collection is bounded by how much garbage a single frame produced, which
for a game loop that reuses its state is modest. There is no unpredictable,
generation-triggered pause of the kind that drops a frame in a Lua game.

The collector's roots are simple, and they are simple \emph{because} the
language is small. The mark phase starts from the constants pool, the globals,
and the VM value stack, and follows the references inside lists and maps.
There are no closures and no upvalues in v0.1, so there are no captured
environments to trace; the object graph is shallow. And because the value
stack is empty between frames, at the moment collection runs the only live
roots are the globals and the constants --- exactly the long-lived state --- so
the mark phase has little to do.

\section{Embedding wick}

wick is meant to be embedded. The public surface is `wick.hpp` alone; embedders
never touch the internal headers. The lifecycle is create, register the
environment, load a program, call its functions each frame, collect at your
safe point, and reset for hot reload.

\begin{lstlisting}[language=C++,basicstyle=\ttfamily\small,keywordstyle=\color{wkKeyword}]
#include "wick.hpp"

wick::VM* vm = wick::create();
std::string err;

// typed natives: a signature DSL the compiler enforces
wick::addNative(vm, "game", "spawn(num, num, str): num",
    [](wick::VM& vm, const wick::Value* a, int) {
        int id = mySpawn(a[0].d, a[1].d, wick::getStr(a[2]));
        return wick::Value::num(id);
    }, err);
wick::addConst(vm, "game", "MAX", 64);

wick::load(vm, source, "main.wick", err);   // compile + run top level
wick::call(vm, "update", dt, true, err);    // each frame
wick::collect(vm);                          // at YOUR safe point
wick::reset(vm);                            // hot-reload support
\end{lstlisting}

\paragraph{The signature DSL.} A native is registered with a signature string
that the compiler parses once and then enforces at every call site. The
grammar is:

\begin{quote}
\ttfamily name(type[, type][, type = default]) [: type]
\end{quote}

\noindent
where a type is `num`, `bool`, `str`, or `list`, with a trailing `?` for an
optional. This one string is how `lt.load_save(str): str?` teaches the
compiler that loading a save returns an optional --- the entire optional
discipline of Chapter~4, for an engine call, is carried by two characters in a
signature. Defaults, which must be numeric literals, make trailing parameters
optional at the call site; the compiler fills them in, so the native always
receives full arity and never has to check `argc`. Registering a native under
a namespace (`"lt"`, `"game"`) qualifies its name; passing a null namespace
puts it in the global scope, which is how the built-ins like `floor` and
`sqrt` are installed.

\paragraph{Reporting errors from a native.} A native signals failure by
calling `wick::setError(vm, msg)`. The VM turns that into a `file:line`
runtime error pointing at the call site, so a native's failure is reported the
same way as any other runtime error and lands on the same error screen. The
value helpers --- `makeStr`, `getStr`, `listLen`, `listGet` --- cover reading
and building the heap values a native receives and returns.

\paragraph{Hot reload.} `wick::reset(vm)` clears all program state --- globals,
functions, the compiled protos --- while leaving the registered environment
(the natives and constants) intact. That is exactly what hot reload needs:
tear down the program, keep the bindings, load the new source. Because
`reset` frees the old program's resources before the new one loads, reloading
never leaks --- the property the tutorial promised.

None of this is large. The embed API is about 120 lines, and the whole
implementation is around two thousand. That smallness is the recurring theme
of this book, and it is nowhere more visible than here: a statically typed,
optional-checked, deterministic scripting language, with a garbage collector
and a typed foreign-function interface, that you can read in an afternoon and
own completely.
